\begin{python0}
from solutions import *; clear()
\end{python0}
\sectionthree{Sums again: Area computation}

I'm going to use the \texttt{for}-loop to compute a certain sum.
This time it's going to be the area under a curve. The computation of areas is extremely important in many areas of science.

The idea is actually very simple: you need to know how to compute the area of rectangles and you need to be able to add lots of these areas.

The idea is pretty simple and illustrates the power of being able to do perform repetitions quickly.

First you draw the graph of the function, for example, y = f(x) = x*x, the standard parabola):

%%\includegraphics[width=4.5in,height=2.7736in]{Pictures/10000000000002B3000001AAC5D1A577C30A5997.png}

Suppose you want to compute the area under this curve and above the
x-axis and from x = 2 to x = 7:

%\includegraphics[width=4.5in,height=2.5165in]{Pictures/10000000000002CD000001911B1B47C442016F8D.png}

You can make an approximation by computing the areas of the rectangles
with
\begin{align*}
\text{x=2 to x=3 with height f(2)}\\
\text{x=3 to x=4 with height f(3)}\\
\text{x=4 to x=5 with height f(4)}\\
\text{x=5 to x=6 with height f(5)}\\
\text{x=6 to x=7 with height f(6)}\\
\end{align*}
There are altogether 5 rectangles. Draw these rectangles into the graph
above.

%\includegraphics[width=4.5in,height=2.5362in]{Pictures/10000000000002CB00000193C61596A9DE7C6C44.png}

Note that for each rectangle, the base has width 1. Also, for each
rectangle, I am using the left endpoint on the width to determine the
height. The areas of these rectangles is smaller than the area under the
parabola from x = 2 to 2 = 7.

\begin{ex}
Write a for-loop to compute the sum of the above
rectangles. Write down the sum of areas. We'll need it
for comparison later.
\end{ex}
Now we \lq\lq move'' the sum of areas of our rectangles towards the parabola
in the following way: we will use \textbf{more} rectangles. Previously the
base of the rectangles had a width of 1. Now I'm going
to use rectangles with base of length 0.5. Altogether we now have 10
rectangles:
\begin{align*}
\text{x=2.0 to x=2.5 with height f(2.0)}\\
\text{x=2.5 to x=3.0 with height f(2.5)}\\
\text{x=3.0 to x=3.5 with height f(3.0)}\\
\text{x=3.5 to x=4.0 with height f(3.5)}\\
\text{x=4.0 to x=4.5 with height f(4.0)}\\
\text{x=4.5 to x=5.0 with height f(4.5)}\\
\text{x=5.0 to x=5.5 with height f(5.0)}\\
\text{x=5.5 to x=6.0 with height f(5.5)}\\
\text{x=6.0 to x=6.5 with height f(6.0)}\\
\text{x=6.5 to x=7.0 with height f(6.5)}\\
\end{align*}
Draw these 10 rectangles into the above graph.

%\includegraphics[width=4.5in,height=2.5492in]{Pictures/10000000000002C9000001940CA977BFEBA1D44B.png}

Do you see that the total area of these rectangles are closer to the
actual area under the parabola than the one from our first
approximation?

Do you see that, intuitively, the total area of the rectangles gets
closer and closer to the area under the parabola from x = 2 to x = 7 as
the number of rectangles used increases? Here's the case
where the length of the base of the rectangles is 0.25:

%\includegraphics[width=4.5in,height=2.5181in]{Pictures/10000000000002C90000018F19E7A9D9A73C9FA9.png}
and here's the case where the length is 0.125

%\includegraphics[width=4.4272in,height=2.139in]{Pictures/10000000000002CA00000191D2EF8AEA529C84AC.png}

Using the code that computes equally spaced point, the sum of the areas
of these rectangles can thus be computed as follows:
\begin{console}
double a = 2;
double b = 7;
int n = 11; // 10 rects mean 11 points
double sum = 0;
double x = 0;
double d = double(b - a) / (n - 1);
for (int i = 0; i < n - 1; ++i)
{   
    sum += d * (x * x);
    x += d;
}

// Why do we NOT need to execute the following?
// x = b;
// sum += d * (x * x);
\end{console}
\begin{ex}
Compute area when 20 rectangles are used.
\end{ex}
\begin{ex}
Now do the same where you use 100 rectangles.
\end{ex}

Now write a program that prompts the user for the number of rectangles!
Say we call this variable n. Run the program with n = 1000 rectangles!!!
You bet the value computed by your program is extremely close to the
real area of the parabola. Of course you wouldn't add
the areas of 1000 rectangles by hand -- it would take too long!
That's what a program is for -- if you know C++!!!

\begin{ex}
If you know some Calculus 1, you know that the area
under the parabola from x = 2 to x = 7 is in fact exactly
\[(1 / 3)(7 * 7 * 7 -- 2 * 2 * 2)\]

Use your calculator (or C++) to compute this value and compare against
the value obtained by your program with n = 1000 rectangles.
\end{ex}

\begin{ex}
The above prompts the user for n, the number of
rectangles to use for the area computation using a for-loop. Write a
nested for-loop where the outer for-loop supplies values of n running
from 1000 to 1000000 by increments of 1000. The pseudocode looks like
this:
\begin{align*}
\text{for n = 1000, 2000, 3000, \ldots, 1000000:}\\
\tab[3em]{\text{compute and print the area using n rectangles}}
\end{align*}
\end{ex}

\begin{ex}
 Modify the above program so that instead of computing
the area under the parabola from x = 2 to x = 7, the program prompts the
user for a, b, n and computes the area under the parabola from x = a to
x = b using n rectangles.
\end{ex}

\begin{ex}
 What is the area under the curve y = x*x*x from x=0
to x=10?
\end{ex}

The computation of areas under a curve is important. For instance in
Physics, the work done by a force F is the area under the
force-displacement graph where you plot the force acting on an object
and the displacement of the object. By Newton's law of
universal gravitation, we know the gravitational force between two
objects which are apart from each other at a fixed distance. So if a
rocket wants to move away from earth, you can compute the work done
moving the rocket to a distance from the surface of earth to a far away
point in space. Knowing the work to be done moving the rocket from earth
through a required distance, you can then compute the amount of chemical
fuel that is needed to combust in order to convert the chemical energy
into required work done. Of course there are many other applications of
the computation of areas.

